\section{Introduction}
Frequent itemset mining (FIM) finds the sets of items that occur together in at least
$\sigma_{\min}$ transactions of a database, where $\sigma_{\min}$ is the minimum support
threshold (called minsup when written as a fraction of the database). It is the first
step of association rule mining and underlies many other pattern-mining tasks. Because
the number of frequent itemsets can grow very quickly as the threshold drops, the way an
algorithm lays out the data and counts support decides how low a threshold it can reach.

Eclat~\cite{zaki2000} stores the database vertically, pairing each item with its tidset,
the set of identifiers of the transactions that contain it. The support of an itemset is
then the size of the intersection of its items' tidsets, and Eclat searches depth-first
over prefix-based equivalence classes. We reimplemented the algorithm from the original
paper in Julia using only the standard library, compared it with SPMF's Eclat and
PrePost+~\cite{spmf2016}, and used it in two applications. The rest of the paper is
organized in three parts:
\begin{itemize}
\item Section~\ref{sec:implementation} presents a baseline that follows the paper and an
optimized variant (bitsets, a search order by increasing support, perfect-extension
pruning, pair counting, and tid lists for sparse data), together with a correctness
argument, the test suite, and the measured gain of each optimization.
\item Section~\ref{sec:evaluation} describes the shared measurement protocol
(Section~\ref{sec:protocol}), checks every itemset against SPMF
(Section~\ref{sec:exact}), and reports runtime against minsup
(Section~\ref{sec:minsup}), the number of frequent itemsets against minsup
(Section~\ref{sec:output-size}), peak memory (Section~\ref{sec:memory}), scalability
with database size (Section~\ref{sec:scalability}), the effect of transaction length
(Section~\ref{sec:width}), and two directions for further optimization
(Section~\ref{sec:future}).
\item Section~\ref{sec:applications} applies Eclat to market basket analysis on Online
Retail II (Section~\ref{sec:basket}) and to image-pattern mining on Semeion
(Section~\ref{sec:visual-patterns}).
\end{itemize}

We follow two rules when presenting results. First, a performance comparison only uses
runs whose output, both itemsets and supports, matches an independent run of SPMF's
Eclat exactly, since two programs can report the same number of itemsets with different
supports. Second, every headline number is tied to the machine and the source revision
that produced it, and every random step uses seed 42. The project repository contains the
source code, configurations, data hashes, and reproduction commands for every experiment.
